% Front page (guide): three approaches, each telling the whole story on one page.
\documentclass[10pt,border=2mm]{standalone}
\usepackage[T1]{fontenc}
\usepackage{figurestyle}
\input{primitives.tex}
\input{molecules.tex}
\input{numbers.tex}
\begin{document}
\begin{tikzpicture}[plate,head/.style={font=\fontsize{9}{11}\selectfont\bfseries,inner sep=0pt},rng/.style={font=\fontsize{8}{10}\selectfont,inner sep=0pt}]
\path[use as bounding box] (0,0) rectangle (15,30);
% ================================================= (A) nested outlines read outside in: the carving as four nested regions
\node[lab,anchor=west] at (.2,29.6) {(A) nested regions, outside in: the whole space, the feasible solid, the version cells, the states};
\begin{scope}[shift={(.5,20.0)}]
  \draw[heavy,rounded corners=10pt] (0,0) rectangle (14,9.0); \node[head,anchor=north west] at (.3,8.75) {Configurations and States}; \node[rng,anchor=north east] at (13.7,8.75) {\S\S\,1--3};
  \path[cut] (.35,.35) rectangle (13.65,6.9); \path[fill=white,rounded corners=9pt] (.9,.5) rectangle (13.1,6.75);
  \draw[heavy,rounded corners=9pt] (.9,.5) rectangle (13.1,6.75); \node[head,anchor=north west] at (1.2,6.5) {Fragments and Feasible Regimes}; \node[rng,anchor=north east] at (12.8,6.5) {\S\S\,4--6};
  \path[cut] (1.25,.85) rectangle (12.75,4.8); \path[fill=white,rounded corners=8pt] (1.8,1.0) rectangle (12.2,4.65);
  \draw[heavy,rounded corners=8pt] (1.8,1.0) rectangle (12.2,4.65); \node[head,anchor=north west] at (2.1,4.4) {Symmetry and Localization}; \node[rng,anchor=north east] at (11.9,4.4) {\S\S\,7--9};
  \path[cut] (2.15,1.35) rectangle (11.85,2.9); \path[fill=white,rounded corners=7pt] (2.7,1.5) rectangle (11.3,2.75);
  \draw[heavy,rounded corners=7pt] (2.7,1.5) rectangle (11.3,2.75); \node[head,anchor=north west] at (3.0,2.5) {Representations and Global Structure}; \node[rng,anchor=north east] at (11.0,2.5) {\S\S\,10--12};
  % emblems: one per ring, placed in the ring's own band
  \begin{scope}\molsolid\mollabelsfalse\def\radH{.16}\def\radX{.25}\pic[scale=.6] at (6.2,7.75) {mol3d-water-X}; \pic[scale=.6] at (8.0,7.75) {mol3d-water-RX};\end{scope}
  \draw[harrow] (6.9,7.75)--(7.3,7.75);
  \begin{scope}[shift={(7.0,5.6)}]\tikzset{tube/Ro=.9,tube/Ri=.62,tube/H=.22,tube/E=.42,tube/a1=-62,tube/a2=8}\pic (ft) at (0,0) {thicktube}; \foreach \p in {top0,top120,top240} {\hatomat{ft-\p}{.05}{ColGold}}\end{scope}
  \begin{scope}[shift={(7.0,3.55)}]\foreach \a in {90,210,330} {\path[fill=ColGold] ({.42*cos(\a)},{.42*sin(\a)}) circle[radius=.16]; \draw[fine] ({.42*cos(\a)},{.42*sin(\a)}) circle[radius=.16];} \draw[fine] (0,0) circle[radius=.42];\end{scope}
  \begin{scope}[shift={(7.0,2.12)}]\foreach \k in {0,1,2} {\path[fill=white] (-.6,{-.2+.15*\k})--(.4,{-.2+.15*\k})--(.6,{-.08+.15*\k})--(-.4,{-.08+.15*\k})--cycle; \draw[fine] (-.6,{-.2+.15*\k})--(.4,{-.2+.15*\k})--(.6,{-.08+.15*\k})--(-.4,{-.08+.15*\k})--cycle;}\end{scope}
\end{scope}
\node[sub,anchor=north west,text width=14cm] at (.5,19.6) {hatched bands are what each step removes: the constraint named by the part; the retained region is carved further inside};
% ================================================= (B) one carved object: C as a slab, F as the tube carved from it, the versions on it, the sheets above
\node[lab,anchor=west] at (.2,18.4) {(B) one object carved through the four parts: slab, tube, versions, sheets; each part is a stage of the same drawing};
\begin{scope}[shift={(.5,9.6)}]
  % the slab C (hatched: everything cut away), the tube F standing in it, the band with versions, two sheets lifted above, the water glyphs at the far left as "a configuration"
  \path[cut] (0,0)--(12.5,0)--(14,1.3)--(1.5,1.3)--cycle; \draw[heavy] (0,0)--(12.5,0)--(14,1.3)--(1.5,1.3)--cycle;
  \node[tag,anchor=north west] at (.2,1.2) {$\conf$};
  \begin{scope}[shift={(7.2,1.9)}]\tikzset{tube/Ro=2.2,tube/Ri=1.55,tube/H=.5,tube/E=.42,tube/a1=-62,tube/a2=8}
    \path[fill=white] (0,-.5) ellipse[x radius=2.2,y radius=.92];
    \pic (bt) at (0,0) {thicktube};
    \draw[tpath] plot[domain=0:115,samples=30,variable=\a] ({1.875*sin(\a)},{.5-.42*1.875*cos(\a)});
    \draw[upath] plot[domain=0:1,samples=30,variable=\s] ({1.875*sin(-60*\s)},{.5-2*.5*\s-.42*1.875*cos(60*\s)});
    \foreach \p in {top0,top120,top240,bot300} {\hatomat{bt-\p}{.08}{ColGold}}
  \end{scope}
  % sheets above the tube: two oblique plates over the version H, lifted
  \foreach \k in {0,1} {\pgfmathsetmacro{\yy}{4.6+.5*\k}\path[fill=white] (5.6,\yy)--(8.2,\yy)--(8.9,\yy+.45)--(6.3,\yy+.45)--cycle; \draw[fine] (5.6,\yy)--(8.2,\yy)--(8.9,\yy+.45)--(6.3,\yy+.45)--cycle;}
  \draw[fine,dash pattern=on 1.6pt off 1.3pt] (7.2,3.1)--(7.2,4.5);
  \begin{scope}\molsolid\mollabelsfalse\pic[scale=.6] at (2.0,3.6) {mol3d-water-X};\end{scope}
  \draw[leader] (2.0,3.2)--(3.0,1.2); \node[sub,anchor=east,text width=2.6cm,align=right] at (1.9,4.4) {a configuration is a point of $\conf$ (\S\S\,1--3)};
  \node[sub,anchor=west,text width=3.0cm] at (10.2,3.2) {$\mathcal F$ carved from $\conf$, its versions on the band (\S\S\,4--6)};
  \node[sub,anchor=west,text width=3.0cm] at (9.3,5.3) {species, spin, packets on the versions (\S\S\,7--9)};
  \node[sub,anchor=east,text width=3.2cm,align=right] at (5.4,5.3) {charts, sheets and modes above (\S\S\,10--12)};
\end{scope}
% ================================================= (C) a strip of four stages, left to right, without frames
\node[lab,anchor=west] at (.2,8.2) {(C) four stages in a strip, left to right, no frames: whole, carved solid, versions, sheets};
\begin{scope}[shift={(.5,2.2)}]
  \foreach \x/\t/\r in {1.6/{Configurations and States}/{1--3},5.3/{Fragments and Feasible Regimes}/{4--6},9.0/{Symmetry and Localization}/{7--9},12.7/{Representations and Global Structure}/{10--12}} {
    \node[head,anchor=north,text width=3.3cm,align=center] at (\x,5.2) {\t}; \node[rng,anchor=north] at (\x,4.4) {\S\S\,\r};}
  \begin{scope}\molsolid\mollabelsfalse\pic[scale=.7] at (1.6,2.6) {mol3d-water-X};\end{scope}
  \begin{scope}[shift={(5.3,2.6)}]\tikzset{tube/Ro=1.1,tube/Ri=.78,tube/H=.27,tube/E=.42,tube/a1=-62,tube/a2=8}\pic (ct) at (0,0) {thicktube}; \foreach \p in {top0,top120,top240} {\hatomat{ct-\p}{.06}{ColGold}}\end{scope}
  \begin{scope}[shift={(9.0,2.6)}]\foreach \a in {90,210,330} {\path[fill=ColGold] ({.6*cos(\a)},{.6*sin(\a)}) circle[radius=.22]; \draw[fine] ({.6*cos(\a)},{.6*sin(\a)}) circle[radius=.22];} \draw[fine] (0,0) circle[radius=.6];\end{scope}
  \begin{scope}[shift={(12.7,2.2)}]\foreach \k in {0,1,2} {\path[fill=white] (-.9,{.25*\k})--(.6,{.25*\k})--(.9,{.2+.25*\k})--(-.6,{.2+.25*\k})--cycle; \draw[fine] (-.9,{.25*\k})--(.6,{.25*\k})--(.9,{.2+.25*\k})--(-.6,{.2+.25*\k})--cycle;}\end{scope}
  \foreach \x in {3.45,7.15,10.85} \draw[harrow] (\x-.3,2.6)--(\x+.3,2.6);
\end{scope}
\node[sub,anchor=north west,text width=14cm] at (.5,1.4) {Critique in docs/illustration-notes.md after rendering.};
\end{tikzpicture}
\end{document}
